%% paper4.tex — Development System: суть, ядро, критические решения
%% Компиляция: lualatex paper4.tex
\documentclass[11pt,a4paper]{article}

\usepackage{fontspec}
\setmainfont{DejaVu Serif}
\setmonofont[Scale=0.88]{DejaVu Sans Mono}

\usepackage{geometry}
\geometry{a4paper, top=28mm, bottom=28mm, left=30mm, right=30mm}

\usepackage{polyglossia}
\setmainlanguage{russian}
\setotherlanguage{english}

\usepackage{microtype}
\usepackage{parskip}
\usepackage{xcolor}
\usepackage{titlesec}
\usepackage{fancyhdr}
\usepackage{hyperref}
\usepackage{float}
\usepackage{enumitem}
\usepackage{booktabs}

\usepackage{listings}

\definecolor{bg}{HTML}{F5F5F2}
\definecolor{kw}{HTML}{00008B}
\definecolor{cm}{HTML}{3C763D}
\definecolor{st}{HTML}{A31515}

\lstdefinelanguage{CL}{
  morekeywords={defun,defpackage,defparameter,let,let*,labels,lambda,
                cond,when,unless,if,progn,handler-case,ignore-errors,
                loop,dolist,mapc,mapcar,reduce,remove-if,format,load,
                in-package,export,setf,push,values,multiple-value-bind,
                append,list,cons,first,rest,car,cdr,null,with-open-file,
                string=,string-upcase,find-symbol,find-package,fboundp,
                probe-file,merge-pathnames,funcall,apply,error,read,
                compile-file,assoc,or,and},
  sensitive=false,
  morecomment=[l]{;},
  morecomment=[s]{\#|}{|\#},
  morestring=[b]",
}

\lstset{
  language=CL,
  backgroundcolor=\color{bg},
  basicstyle=\small\ttfamily,
  keywordstyle=\color{kw}\bfseries,
  commentstyle=\color{cm}\itshape,
  stringstyle=\color{st},
  frame=single, framesep=5pt, rulecolor=\color{gray!20},
  showstringspaces=false,
  breaklines=true, tabsize=2,
  numbers=none,
  xleftmargin=10pt, xrightmargin=4pt,
  captionpos=b,
}

\hypersetup{colorlinks=true, linkcolor=black, urlcolor=blue!60!black}
\titleformat{\section}{\large\bfseries}{\thesection.}{0.6em}{}
\titleformat{\subsection}{\normalsize\bfseries}{\thesubsection.}{0.5em}{}

\pagestyle{fancy}
\fancyhf{}
\fancyhead[L]{\small\textit{Development System — суть и ядро}}
\fancyhead[R]{\small\thepage}
\renewcommand{\headrulewidth}{0.4pt}

\begin{document}

\begin{titlepage}
\centering\vspace*{3cm}
{\LARGE\bfseries Development System\\[0.5em]
\large Суть, ядро, критические решения}

\vspace{1.5cm}
{\large Артемий Уродовских}\\[0.2em]
{\normalsize Март 2026}

\vspace{2.5cm}
\begin{abstract}
\noindent
Эта статья не про Telegram и не про HTTP. Она про то, как система организована внутри: как два независимых процесса договариваются о том, что выполнять; почему поток — это один файл с одной функцией; что значит «персистентный образ» и почему это не просто оптимизация; и как система замыкает петлю, порождая собственные модули через LLM.

Центральные темы: протокол активации, явное состояние в SICP-стиле, образ как реестр, метациркулярность.
\end{abstract}
\end{titlepage}

\tableofcontents
\newpage

% ═══════════════════════════════════════════════════════════════════════════════
\section{Центральная идея}

Система состоит из двух SBCL-процессов. \textbf{Мастер} — управляющий: принимает команды, управляет потоками, хранит состояние диалога. \textbf{Подмастерье} — исполняющий: держит живой образ, загружает потоки в него, вызывает напрямую.

Между ними нет сокета, нет очереди, нет RPC. Они общаются через файловую систему — один файл с именем \texttt{активный.scm}.

Это намеренно простое решение, и в нём есть своя элегантность: каждый процесс можно перезапустить независимо, не разрушив состояние другого. Протокол восстанавливается при следующем чтении файла.

% ═══════════════════════════════════════════════════════════════════════════════
\section{Протокол активации}
\label{sec:activation}

\subsection{Файл как IPC}

Мастер сигнализирует Подмастерью какой поток должен работать, записывая S-выражение:

\begin{lstlisting}[caption={мастер/потоки/активный.scm}]
(активный
  (файл  "/srv/development_system/мастер/потоки/уточнение.lisp")
  (пакет "поток-уточнение"))
\end{lstlisting}

Подмастерье читает это при \textit{каждом} входящем сообщении:

\begin{lstlisting}[caption={подмастерье/main.lisp}]
(defun прочитать-активный ()
  (let ((путь (format nil "~aактивный.scm" *каталог-потоков*)))
    (handler-case
        (with-open-file (f путь)
          (let* ((д     (read f nil nil))
                 (файл  (cadr (assoc 'файл  (cdr д))))
                 (пакет (cadr (assoc 'пакет (cdr д)))))
            (when (and файл пакет) (list файл пакет))))
      (error () nil))))
\end{lstlisting}

Заметим: \texttt{read} читает S-выражение напрямую — никакого парсера, никакого формата. Файл \textit{уже} является структурой данных CL. Это следствие гомоиконности: данные и код имеют одно представление.

\subsection{Семантика «eventual consistency»}

Подмастерье не подписывается на изменения файла. Нет inotify, нет polling на файл. Просто: при каждом запросе — читай файл заново.

Это означает, что смена активного потока вступает в силу на следующем сообщении пользователя, не немедленно. Для интерактивного бота это приемлемо, и это устраняет целый класс проблем с синхронизацией.

\subsection{Что происходит при переключении потока}

\begin{enumerate}[noitemsep]
  \item Мастер пишет новый \texttt{активный.scm}
  \item Пользователь отправляет следующее сообщение
  \item Подмастерье читает файл, получает путь и имя пакета
  \item Проверяет реестр: загружен ли этот пакет?
  \item Если нет — \texttt{load} файла в живой образ, запись в реестр
  \item \texttt{funcall} на \texttt{ВЫПОЛНИТЬ} в загруженном пакете
\end{enumerate}

Шаги 4–5 идемпотентны: один и тот же поток загружается не более одного раза за жизнь процесса.

% ═══════════════════════════════════════════════════════════════════════════════
\section{Персистентный образ}
\label{sec:image}

\subsection{Что это значит}

SBCL (как и любой Common Lisp) — это не интерпретатор, запускающий скрипты. Это \textit{живой образ}: объекты, пакеты, скомпилированные функции живут в памяти процесса до его завершения. \texttt{load} добавляет новый код в этот образ, не перезапуская процесс.

Когда Подмастерье загружает \texttt{уточнение.lisp}, пакет \texttt{:поток-уточнение} становится частью образа навсегда (до перезапуска). Второй вызов того же потока — это просто:

\begin{lstlisting}
(funcall #'поток-уточнение:выполнить задача)
\end{lstlisting}

Никакого файлового ввода-вывода, никакого subprocess, никакой сериализации.

\subsection{Образ как реестр}

Вместо явного реестра-хэша Подмастерье использует сам образ как истину о том, что загружено. Проверка тривиальна:

\begin{lstlisting}[caption={Образ как реестр — find-package как ключ}]
(defun загрузить-в-образ (путь пакет реестр)
  (if (assoc пакет реестр :test #'string=)
      реестр                        ; уже есть — ничего не делать
      (handler-case
          (progn
            (load путь :verbose nil :print nil)
            (cons (cons пакет путь) реестр))  ; добавить в реестр
        (error (e)
          (format t "Ошибка загрузки ~a: ~a~%" пакет e)
          реестр))))
\end{lstlisting}

Реестр — \texttt{alist} вида \texttt{("поток-уточнение" . "/path/уточнение.lisp")}. Он не хранится в глобальной переменной — он \textit{живёт как параметр рекурсии} (подробнее в следующем разделе).

\subsection{Что происходит при ошибке загрузки}

Ошибка компиляции или выполнения при \texttt{load} ловится \texttt{handler-case} и реестр возвращается без изменений. Поток считается «не загруженным», и при следующей попытке загрузка будет повторена. Это консервативное поведение: лучше попробовать снова, чем навсегда записать неработающий поток.

% ═══════════════════════════════════════════════════════════════════════════════
\section{Явное состояние: SICP в действии}
\label{sec:state}

\subsection{Проблема}

Цикл обработки сообщений имеет состояние: реестр загруженных потоков и текущий offset Telegram. Самый простой способ — хранить их в глобальных переменных. Это работает, но создаёт скрытые зависимости: любая функция может случайно мутировать реестр, порядок операций имеет значение, тестирование затруднено.

\subsection{Решение}

Состояние передаётся явно как параметр рекурсии. Главный цикл Подмастерья:

\begin{lstlisting}[caption={подмастерье/main.lisp — рекурсия с явным состоянием}]
(defun опрос (сдвиг реестр)
  (handler-case
      (let* ((сообщения (извлечь-сообщения
                          (разобрать-обновления
                            (получить-обновления сдвиг)))))
        (multiple-value-bind (новый-рег новый-сдвиг)
            (обр-сообщения сообщения реестр сдвиг)
          (when (> новый-сдвиг сдвиг)
            (сохранить-сдвиг новый-сдвиг))
          (опрос новый-сдвиг новый-рег)))
    (error (e)
      (format t "Ошибка цикла: ~a~%" e)
      (sleep 5)
      (опрос сдвиг реестр))))  ; перезапуск с тем же состоянием
\end{lstlisting}

Аргументы \texttt{сдвиг} и \texttt{реестр} — это всё состояние системы. Нет глобальных переменных, нет скрытых зависимостей.

\subsection{Обработка списка сообщений}

Каждая порция сообщений обрабатывается рекурсивно:

\begin{lstlisting}[caption={подмастерье/main.lisp — обработка с накоплением}]
(defun обр-сообщения (список реестр макс)
  (if (null список)
      (values реестр макс)
      (let* ((с       (car список))
             (ид-чата (second с))
             (текст   (third с)))
        (multiple-value-bind (новый-рег)
            (if (админ? ид-чата)
                (multiple-value-bind (ответ р)
                    (обработать текст ид-чата реестр)
                  (послать-сообщение ид-чата ответ)
                  р)
                реестр)
          (обр-сообщения (cdr список) новый-рег
                          (max макс (1+ (first с))))))))
\end{lstlisting}

Паттерн из SICP главы 3: вместо мутабельного накопителя — рекурсия с аккумулятором. Реестр обновляется по мере загрузки потоков; каждая итерация получает свежую копию.

\subsection{Что даёт явное состояние}

\begin{itemize}[noitemsep]
  \item При ошибке цикла \texttt{опрос} перезапускается с \textit{последним корректным} состоянием — реестром и сдвигом до момента сбоя.
  \item Тестирование: можно вызвать \texttt{обр-сообщения} с любым реестром, не меняя глобального состояния.
  \item Нет гонки: в однопоточном процессе явное состояние полностью предсказуемо.
\end{itemize}

% ═══════════════════════════════════════════════════════════════════════════════
\section{Контракт потока}
\label{sec:contract}

\subsection{Почему одна функция}

Каждый поток определяет ровно одну точку входа:

\begin{lstlisting}
(defpackage :поток-ИМЯ (:use :cl) (:export #:выполнить))
(in-package :поток-ИМЯ)
(defun выполнить (задача) ...)  ; строка → строка
\end{lstlisting}

Это минимальный контракт. Он достаточен потому что:

\begin{enumerate}[noitemsep]
  \item Обнаружение через рефлексию: \texttt{find-package} + \texttt{find-symbol} + \texttt{fboundp} — три проверки, не требующих метаданных.
  \item LLM-генерация: промпт описывает контракт в трёх строках, и модель его соблюдает.
  \item Тестирование в REPL: \texttt{(поток-эхо:выполнить "тест")} — ничего лишнего.
\end{enumerate}

\subsection{Поиск точки входа через рефлексию}

\begin{lstlisting}[caption={мастер/cl/потоки.lisp — рефлексия без метаданных}]
(defun %pkg-name (имя)
  (string-upcase (concatenate 'string "поток-" имя)))

(defun %найти-выполнить (имя)
  (let* ((pkg (find-package (%pkg-name имя)))
         (sym (when pkg (find-symbol "ВЫПОЛНИТЬ" pkg))))
    (when (and sym (fboundp sym)) (symbol-function sym))))
\end{lstlisting}

Этот код не знает ничего о содержимом потока. Он только спрашивает: «есть ли в образе пакет с таким именем, и определена ли в нём функция ВЫПОЛНИТЬ?» Если да — вызывает. Иначе — загружает файл и пробует снова.

% ═══════════════════════════════════════════════════════════════════════════════
\section{Метациркулярность: поток порождает потоки}
\label{sec:meta}

\subsection{Структура}

Порождатель — поток, соответствующий тому же контракту (\texttt{выполнить/1}), но его реализация:

\begin{lstlisting}[caption={мастер/потоки/порождатель.lisp — суть}]
(defun выполнить (описание)
  (handler-case
      (let* ((код   (извлечь-из-блока
                      (запросить-генерацию описание)))
             (имя   (имя-пакета код))
             (путь  (записать-поток имя код))
             (ok?   (компилировать путь)))
        (when ok? (загрузить-в-образ путь
                    (format nil "поток-~a" имя)))
        (format nil "Поток «~a» ~a." имя
                (if ok? "готов" "не скомпилирован")))
    (error (e) (format nil "Ошибка: ~a" e))))
\end{lstlisting}

Цепочка: \textit{текст} → LLM → \textit{код} → \texttt{compile-file} → \texttt{load} → \textit{живая функция}.

\subsection{Почему это работает без перезапуска}

\texttt{compile-file} компилирует \texttt{.lisp} в \texttt{.fasl}. \texttt{load} загружает \texttt{.fasl} в текущий образ. После этого \texttt{find-package} находит новый пакет, и следующий вызов \texttt{\%найти-выполнить} возвращает свежесозданную функцию.

Процесс не знает заранее о новом потоке. Образ обнаруживает его через рефлексию. Это не отличается от любого другого потока — механизм универсален.

\subsection{Извлечение имени пакета из кода}

\begin{lstlisting}[caption={мастер/потоки/порождатель.lisp — имя без парсера}]
(defun имя-пакета (код)
  (let* ((метка "поток-")
         (поз   (search метка код)))
    (if поз
        (let* ((нач (+ поз (length метка)))
               (кон (position-if
                      (lambda (c) (member c '(#\Space #\Newline #\( #\))))
                      код :start нач)))
          (subseq код нач (or кон (length код))))
        "безымянный")))
\end{lstlisting}

Не парсер CL, а простой текстовый поиск. Это работает потому что промпт жёстко требует \texttt{:поток-ИМЯ} в начале файла. Хрупко, но достаточно.

\subsection{Замкнутая петля}

\begin{figure}[H]\centering\small\begin{verbatim}
Пользователь пишет: /породить переводчик текста

Мастер → Порождатель.выполнить("переводчик текста")
  → запросить-генерацию → LLM
  → код на CL: (defpackage :поток-переводчик ...)
  → записать-поток → мастер/потоки/переводчик.lisp
  → compile-file → переводчик.fasl
  → load → :поток-переводчик теперь в образе

Пользователь пишет: /запустить переводчик Привет мир

Мастер → %найти-выполнить("переводчик")
  → find-package("ПОТОК-ПЕРЕВОДЧИК") → найден
  → funcall #'поток-переводчик:выполнить → "Hello, world"
\end{verbatim}
\caption{Полный цикл порождения и использования}
\end{figure}

Между порождением и первым вызовом нет перезапуска процесса. Образ расширяется в рантайме.

% ═══════════════════════════════════════════════════════════════════════════════
\section{Ключевые архитектурные решения}
\label{sec:decisions}

\subsection{Файл вместо IPC}

Мастер и Подмастерье могли бы общаться через Unix-сокет, pipe или разделяемую память. Вместо этого — файл \texttt{активный.scm}.

Это проигрывает по latency (файловый ввод-вывод), но выигрывает по:
\begin{itemize}[noitemsep]
  \item \textbf{Отлаживаемости}: состояние видно через \texttt{cat активный.scm}
  \item \textbf{Отказоустойчивости}: при перезапуске одного процесса второй не замечает разрыва
  \item \textbf{Простоте}: нет кода установки соединения, нет сериализации, нет reconnect-логики
\end{itemize}

Для интерактивного бота, где latency определяется сетью и LLM, это приемлемый компромисс.

\subsection{S-выражение как формат конфигурации}

И \texttt{активный.scm}, и файл «души» (\texttt{душа.scm}) — это S-выражения, читаемые стандартным \texttt{read}. Не JSON, не YAML, не TOML.

Это означает: редактирование конфигурации — это написание CL-данных. REPL может читать и модифицировать конфигурацию теми же инструментами, что и код. Граница между «данными» и «кодом» размыта намеренно.

\subsection{Однопоточность как инвариант}

Оба процесса однопоточные. Нет мьютексов, нет атомарных операций, нет race conditions внутри процесса. Состояние предсказуемо.

Цена: если поток выполняется долго (LLM-запрос на 10 секунд), бот не реагирует на другие сообщения. Для текущего сценария (один пользователь-администратор) это приемлемо.

\subsection{Процессное разделение как изоляция}

Если поток в Подмастерье падает с необрабатываемой ошибкой — Мастер продолжает работать. Если Мастер падает — Подмастерье держит текущий поток активным. Разделение процессов — это бесплатная изоляция отказов.

\subsection{Образ как единица деплоя}

Подмастерье не перезапускается для обновления потока. Достаточно записать новый \texttt{.lisp} файл и вызвать \texttt{load}. Это принцип Lisp Machine: система эволюционирует без остановки.

% ═══════════════════════════════════════════════════════════════════════════════
\section{Что система знает о потоках и чего не знает}

Это важно для понимания ограничений.

\begin{table}[H]\centering\small
\begin{tabular}{p{5cm}p{5cm}}
\toprule
\textbf{Система знает} & \textbf{Система не знает} \\
\midrule
Имя пакета потока & Что поток делает \\
Путь к файлу & Какие ресурсы использует \\
Факт загрузки в образ & Сколько времени займёт выполнение \\
Наличие функции \texttt{выполнить} & Может ли поток завершиться \\
\bottomrule
\end{tabular}
\caption{Граница знания системы о потоках}
\end{table}

Система намеренно не инспектирует содержимое потока. Это делает её обобщённой: любой код с нужным контрактом становится потоком. Но это также означает, что зависший поток зависает весь Подмастерье. Это известное ограничение.

% ═══════════════════════════════════════════════════════════════════════════════
\section{Об обработке ошибок}

\subsection{Текущий подход}

Система использует \texttt{handler-case} на всех внешних границах (HTTP, файловый ввод-вывод, \texttt{load}). Ошибки логируются, функции возвращают \texttt{nil} или строку-заглушку. Процесс продолжает работу.

Это «defensive programming»: система предпочитает продолжить с деградировавшим поведением, а не упасть.

\subsection{Что теряется}

Common Lisp предоставляет значительно более богатую систему условий: \texttt{restart-case} позволяет сигнализировать условие и предложить \textit{рестарты} — точки восстановления, из которых вызывающий код выбирает нужный.

Пример того, как могло бы быть:

\begin{lstlisting}[caption={Гипотетический API с рестартами}]
;; Вместо: ignore-errors или handler-case с nil
(restart-case
    (load поток-путь)
  (retry ()
    :report "Попробовать загрузить снова"
    (load поток-путь))
  (skip ()
    :report "Пропустить этот поток"
    реестр)
  (use-cached (старый-пакет)
    :report "Использовать предыдущую версию"
    старый-пакет))
\end{lstlisting}

В текущей архитектуре решение «что делать при ошибке» принято заранее внутри функции. В полноценной CL-системе это решение могло бы приниматься вызывающим кодом без размотки стека — именно так работает интерактивный debugger SBCL: он останавливается в точке ошибки и предлагает рестарты.

Текущий подход — компромисс: проще читать, проще отлаживать в продакшене, но теряет информацию и восстанавливаемость.

% ═══════════════════════════════════════════════════════════════════════════════
\section{Сводка}

\begin{table}[H]\centering\small
\begin{tabular}{p{4cm}p{8cm}}
\toprule
\textbf{Решение} & \textbf{Следствие} \\
\midrule
Файл как IPC & Отлаживаемость, независимость процессов, eventual consistency \\
S-выражение как конфиг & Единый язык для кода и данных, REPL-доступ к конфигурации \\
Персистентный образ & Потоки загружаются один раз, прямой \texttt{funcall} без overhead \\
Явное состояние & Предсказуемость, тестируемость, корректный restart после ошибки \\
Один пакет + одна функция & Обнаружение через рефлексию, LLM-генерируемость \\
Однопоточность & Нет race conditions, нет мьютексов \\
\texttt{handler-case} на границах & Продолжение при сбоях, потеря восстанавливаемости \\
\bottomrule
\end{tabular}
\caption{Ключевые решения и их следствия}
\end{table}

Система небольшая — около 1100 строк. Каждое из этих решений можно проследить до конкретного места в коде. Это и есть критерий хорошей архитектуры: между намерением и реализацией нет зазора.

\end{document}
